\section*{\tr{Resumo}{Abstract}}
\addcontentsline{toc}{section}{\tr{Resumo}{Abstract}}

\pt{%
Acidentes em operações de pátio logístico, como colisões de caminhões de manobra com carretas, docas e estruturas, são eventos raros, de alto custo e distribuídos de forma desigual entre veículos e unidades operacionais. Na prática, a gestão de segurança costuma comparar unidades por contagens brutas de acidentes, sem considerar quanto cada unidade operou. Como o volume de operação muda muito entre unidades e ao longo do tempo, esses rankings confundem risco com exposição e oscilam por acaso. Este projeto propõe modelar a frequência de acidentes por veículo e mês com modelos hierárquicos de contagem (Poisson e binomial negativa com efeitos aleatórios), usando como exposição as horas e distâncias registradas pela telemetria embarcada e, como preditores, taxas de eventos de condução registrados pelos sensores, como frenagens bruscas, curvas acentuadas e excesso de velocidade. O projeto também propõe estimar a subnotificação de acidentes, comparando colisões detectadas pelos sensores com os relatos formais. Os modelos serão avaliados fora da amostra, por desempenho preditivo e pela estabilidade dos rankings de risco. Caso a telemetria não possa ser usada, o mesmo arcabouço é aplicado por unidade e mês, com horas trabalhadas como exposição. Como extensão opcional, a previsão probabilística da exposição futura será combinada ao modelo de risco para projetar acidentes esperados e apoiar o planejamento da supervisão de segurança.
}
\en{%
Accidents in logistics yard operations, such as collisions of yard trucks with trailers, docks and structures, are rare, costly events that are unevenly distributed across vehicles and operating sites. In practice, safety management often compares sites by raw accident counts, ignoring how much each site operated. Since operating volume varies widely across sites and over time, such rankings confuse risk with exposure and fluctuate by chance. This project proposes to model accident frequency per vehicle and month with hierarchical count models (Poisson and negative binomial with random effects), using the hours and distances recorded by on-board telematics as exposure and, as predictors, rates of driving events recorded by the sensors, such as harsh braking, harsh turns and speeding. The project also proposes to estimate accident underreporting by comparing collisions detected by the sensors with formal reports. Models will be evaluated out of sample, by predictive performance and by the stability of risk rankings. If telematics cannot be used, the same framework is applied per site and month, with hours worked as exposure. As an optional extension, a probabilistic forecast of future exposure will be combined with the risk model to project expected accidents and support the planning of safety supervision.
}

\medskip
\noindent\textbf{\tr{Palavras-chave}{Keywords}:} \tr%
  {modelos de contagem; modelos hierárquicos; exposição ao risco; telemetria veicular; segurança do trabalho; subnotificação.}%
  {count models; hierarchical models; risk exposure; vehicle telematics; occupational safety; underreporting.}
